\documentclass[11pt]{article}
\usepackage[margin=1in]{geometry}
\usepackage{amsmath,amssymb,amsthm,mathtools,microtype,hyperref}
\hypersetup{hidelinks}
\newtheorem{theorem}{Theorem}
\newtheorem{lemma}[theorem]{Lemma}
\newtheorem{proposition}[theorem]{Proposition}
\newtheorem{corollary}[theorem]{Corollary}
\theoremstyle{remark}\newtheorem{remark}[theorem]{Remark}
\newcommand{\dA}{\mathrm dA}
\newcommand{\ddc}{dd^c}
\newcommand{\CP}{\mathbb P}
\newcommand{\R}{\mathbb R}
\newcommand{\C}{\mathbb C}
\newcommand{\OO}{\mathcal O}
\newcommand{\II}{\mathcal I}
\newcommand{\mm}{\mathfrak m}
\DeclareMathOperator{\Var}{Var}
\title{Anisotropic growing point jets on compact products\\
and transfer by non-descending fiber selectors}
\author{}
\date{4 October 2026}
\begin{document}
\maketitle
\begin{abstract}
We formulate an explicit higher-dimensional growing-jet theorem for compact
products of projective lines. Unequal curvature degeneracies and nonzero smooth
radial remainders are allowed. For every sublinear growing vanishing order,
the actual saddle center and variance give a Gaussian transition. The leading
homogeneous center has a smaller, sharp range in this class. A uniform estimate
for the full compact monomial norms proves the global comparison, including
arbitrarily slowly growing orders. An isolated singularity on a compact
connected-fiber holomorphic submersion selects exactly the base point jets,
although the singularity and its multiplier ideal do not descend. Matched
smooth volumes give exact Hilbert isometries and preserve the intrinsic base
growth exponent. The results are sufficient conditions for a specified radial
class; mixed remainders and arbitrary fiber selectors are outside their scope.
\end{abstract}

\section{Compact data and the transition}
We use $\ddc=(i/(2\pi))\partial\bar\partial$ and Euclidean area $\dA$.
Fix $n\geq2$, integers $r_j\geq2$, with at least two unequal, and
$\varepsilon_j>0$. Set $d_j=r_j+1$ and
\[
 Y=(\CP^1)^n,\quad x=(0,\ldots,0),\quad
 L=\boxtimes_{j=1}^n\OO_{\CP^1}(d_j).
\]
In the product affine frame, with $t_j=|z_j|^2$, put
\begin{align}
 |e|_h^2&=e^{-\phi},&
 \phi(z)&=\sum_j\psi_j(t_j),&
 \psi_j(t)&=\log(1+t^{r_j}+\varepsilon_jt^{r_j+1}),\label{eq:metric}\\
 \nu_Y&=a(t)\prod_j\dA_j,&
 a(t)&=\prod_j a_j(t_j),& a_j(t)&=(1+t)^{-2}.
\end{align}
Define
\[
 H_j(t)=t\psi_j'(t),\qquad G_j(t)=tH_j'(t),\qquad
 g_j(t)=G_j(t)/t=\psi_j'(t)+t\psi_j''(t),\qquad
 \omega=c_1(L,h).
\]
For integers $q=q_m\to\infty$, $q=o(m)$, give
\[
 V_m=H^0(Y,L^m\otimes\mm_x^q),\qquad
 \|s\|^2=\int_Y|s|_{h^m}^2\nu_Y
\]
the indicated smooth norm. Its diagonal kernel is
$B_m=\sum_\ell |s_\ell|_{h^m}^2$ for an orthonormal basis.

\begin{theorem}\label{thm:base}
For every observation sequence with $t_j\to0$,
\[
 \mu_j=mH_j(t_j),\qquad \mu_j/q\longrightarrow\alpha_j>0,
 \qquad \sum_j\alpha_j=1,
\]
set $V=\sum_j mG_j(t_j)$. If
$ (\sum_j\mu_j-q)/\sqrt V\to\tau\in\R$, then
\begin{equation}\label{eq:transition}
 \frac{\pi^n a(t)B_m(z)}{m^n\prod_jg_j(t_j)}\longrightarrow
 \Phi(\tau),\qquad
 \Phi(\tau)=\frac1{\sqrt{2\pi}}\int_{-\infty}^{\tau}e^{-s^2/2}\,ds.
\end{equation}
All choices of coordinate arguments are allowed. The dimension is eventually
\begin{equation}\label{eq:dimension}
 N_m=\prod_j(d_jm+1)-\binom{q+n-1}{n},\qquad
 N_m/m^n\longrightarrow D:=\prod_jd_j.
\end{equation}
In particular the section-growth exponent is $n$. Globally,
\begin{equation}\label{eq:measure}
 m^{-n}B_m\nu_Y\longrightarrow\omega^n/n!
 \quad\text{in total variation}.
\end{equation}
\end{theorem}

An explicit corrected observation scale is obtained by fixing
$\alpha_j>0$, $\sum\alpha_j=1$, and writing
\begin{equation}\label{eq:corrected}
 R=\sum_j\alpha_jr_j,\qquad T_m=q+\tau\sqrt{Rq},\qquad
 t_j=H_j^{-1}(\alpha_jT_m/m).
\end{equation}
Then $\sum\mu_j=T_m$ and $V\sim Rq$, so \eqref{eq:transition} holds
throughout $q=o(m)$. The radii have leading orders
$(\alpha_jq/(r_jm))^{1/(2r_j)}$.

\begin{proposition}\label{prop:shift}
Instead prescribe the homogeneous radii
\begin{equation}\label{eq:homogeneous}
 t_j=\left(\frac{\alpha_j T_m}{r_jm}\right)^{1/r_j},
 \qquad r_*:=\max_j r_j.
\end{equation}
For this radial class the unshifted law
\[
 \frac{\pi^nB_m(z)}{m^n\prod_j r_j^2|z_j|^{2r_j-2}}
 \longrightarrow\Phi(\tau)
\]
holds for every fixed $\tau$ along the whole sequence if and only if
$q=o(m^{2/(r_*+2)})$. If $q/m^{2/(r_*+2)}\to c\in(0,\infty)$,
the limit is
\[
 \Phi\left(\tau+C_\alpha c^{(r_*+2)/(2r_*)}\right),\qquad
 C_\alpha=\frac1{\sqrt R}\sum_{j:r_j=r_*}
 (r_*+1)\varepsilon_j(\alpha_j/r_*)^{1+1/r_*}.
\]
If that growth ratio tends to infinity while $q=o(m)$, the limit is $1$.
\end{proposition}

The hypotheses describe genuine global metrics. In the infinity chart of a
factor, with $u=|w|^2=1/t$, the weight after its frame change is
$\log(\varepsilon_j+u+u^{d_j})$. It is smooth with positive curvature
at $u=0$. On $0<t<\infty$, $H_j$ is the mean and $G_j$ the variance
of the three distinct exponents $0,r_j,d_j$, with probabilities proportional
to $1,t^{r_j},\varepsilon_jt^{d_j}$. Thus $G_j>0$. At the origin
\begin{align}
 \psi_j(t)&=t^{r_j}+\varepsilon_jt^{r_j+1}+O(t^{2r_j}),\label{eq:expansion}\\
 H_j(t)&=r_jt^{r_j}+(r_j+1)\varepsilon_jt^{r_j+1}+O(t^{2r_j}),\\
 G_j(t)&=r_j^2t^{r_j}+(r_j+1)^2\varepsilon_jt^{r_j+1}+O(t^{2r_j}).
\end{align}
The curvature is semipositive; in the product it loses rank along coordinate
divisors. The remainders in \eqref{eq:expansion} are genuinely nonzero.
The volume also extends smoothly and positively at infinity, and $a(0)=1$.

\section{Uniform global moment estimates}
We prove the estimate needed for moving orders rather than appeal to a
fixed-order localization theorem. Suppress the coordinate index and set
\[
 P(t)=1+t^r+\varepsilon t^{r+1},\quad \psi=\log P,\quad
 a(t)=(1+t)^{-2},\quad d=r+1,
\]
where $r\geq2$. Write $H=t\psi'$, $G=tH'$, $g=G/t$, and define
\begin{equation}\label{eq:moment}
 J_{m,k}=\pi\int_0^\infty t^ke^{-m\psi(t)}a(t)\,dt,
 \quad w_{m,k}(t)=\frac{t^ke^{-m\psi(t)}}{J_{m,k}},
 \quad K_m(t)=\sum_{k=0}^{dm}w_{m,k}(t).
\end{equation}
These compact norms are finite even at $k=dm$, since the integrand is
$O(t^{k-dm-2})$ at infinity.

\begin{lemma}\label{lem:moments}
If $t=t_m\to0$ and $\mu=mH(t)\to\infty$, then, with $V=mG(t)$,
\begin{equation}\label{eq:fullkernel}
 K_m(t)=\frac{V}{\pi a(t)t}(1+o(1)),\qquad
 \sum_k\frac{w_{m,k}(t)}{K_m(t)}\,
 \delta_{(k-\mu)/\sqrt V}\ \Longrightarrow\ N(0,1).
\end{equation}
There is no lower bound on how fast $\mu$ grows. For every fixed $t>0$,
locally uniformly on $(0,\infty)$,
$K_m(t)=mg(t)/(\pi a(t))(1+o(1))$.
If $k\to\infty$, $k=o(m)$, and $H(s)=k/m$, then uniformly along all
such sequences
\begin{equation}\label{eq:laplace}
 J_{m,k}=\pi a(s)s\sqrt{\frac{2\pi}{mG(s)}}
 \exp(k\log s-m\psi(s))(1+o(1)).
\end{equation}
\end{lemma}

\begin{proof}
Put $F(u)=\psi(e^u)$ and $A(u)=a(e^u)e^u$. Then $F'=H(e^u)$,
$F''=G(e^u)>0$. Near zero $H\asymp t^r$, $G/H\to r$, and every
fixed positive-order $u$ derivative of $F$ is $O(H)$. Choose $b>0$
small and constants
\[
 1<\rho_-<r<\rho_+,
 \qquad \rho_-\leq G(t)/H(t)\leq\rho_+\quad(0<t\leq b).
\]
The logarithmic derivatives of $A$ are bounded, and $A(u)\asymp e^u$
for $e^u\leq b$. There is a unique saddle $s$ for $0<k<dm$ because
$H$ increases from $0$ to $d$.

Let $v=\log s$ and
$E(h)=m(F(v+h)-F(v))-kh$. Choose $\eta_0>0$ small enough that
$k\leq\eta_0m$ implies $s\leq b/e^2$ and $k<mH(b)/2$.
On $|h|\leq1$ one has
\[
 E(h)\geq ck h^2,\qquad
 E(h)=\tfrac12mG(s)h^2+O(k|h|^3).
\]
Integration of the elasticity bounds gives
\begin{align*}
 E(h)&\geq k\left(|h|-\frac{1-e^{\rho_-h}}{\rho_-}\right)
 &&(h\leq-1),\\
 E(h)&\geq k\left(\frac{e^{\rho_-h}-1}{\rho_-}-h\right)
 &&(1\leq h\leq\log(b/s)).
\end{align*}
The amplitude ratio is at most $Ce^h$. Outside any fixed small saddle
neighborhood these estimates give exponentially small normalized tails
in $k$. At $h_b=\log(b/s)$, after reducing $\eta_0$,
$E(h_b)\geq cm$: the positive term $m\psi(b)$ dominates
$m\psi(s)=O(k)$ and
$k\log(b/s)\leq Cm\eta_0\log(1/\eta_0)$.
For $h\geq h_b$, $E'(h)\geq cm$. Consequently the remaining normalized
tail is bounded by
\[
 C(b/s)e^{-cm}/(cm-1).
\]
Since $s\asymp(k/m)^{1/r}$ for $1\leq k\leq\eta_0m$, this is a
polynomial factor times $e^{-cm}$.
Rescale $h=y/\sqrt{mG(s)}$. Taylor expansion and the preceding domination
give the Gaussian integral and prove \eqref{eq:laplace}. The same local
argument, without taking a limit, gives
\begin{equation}\label{eq:lower-moment}
 J_{m,k}\geq c a(s)s k^{-1/2}e^{k\log s-m\psi(s)}
 \qquad(1\leq k\leq\eta_0m).
\end{equation}
Indeed one may integrate on $|h|\leq k^{-1/2}$.

Now let $t\to0$, $\mu=mH(t)\to\infty$, so $\mu=o(m)$ and
$V\sim r\mu$. For $|k-\mu|\leq M\sqrt V$, the saddle equation gives
\[
 \log(s/t)=\frac{k-\mu}{V}+O((k-\mu)^2/\mu^2).
\]
Define
$I(k)=m(\psi(t)-\psi(s))-k\log(t/s)$, extending $k$ to real
values through its saddle. Then $I(\mu)=I'(\mu)=0$ and
$I''(k)=1/(mG(s))$. Formula \eqref{eq:laplace} yields
\begin{equation}\label{eq:central}
 w_{m,k}(t)=\frac{\sqrt V}{\pi a(t)t\sqrt{2\pi}}
 \exp\left(-\frac{(k-\mu)^2}{2V}\right)(1+o(1))
\end{equation}
uniformly on each such interval. The natural sum scale is
$T=V/(\pi a(t)t)$.

We give summable estimates for the other indices. For
$1\leq k\leq\eta_0m$, \eqref{eq:lower-moment} implies
\[
 \frac{w_{m,k}}T\leq C\frac{\sqrt k}{\mu}\frac ts e^{-I(k)}.
\]
Along the interval between $k$ and $\mu$,
$I''(k)\geq1/(\rho_+k)$, so
\[
 I(k)\geq\rho_+^{-1}\bigl(k\log(k/\mu)-k+\mu\bigr).
\]
If $k\leq\mu$, then $t/s\leq(\mu/k)^\alpha$ with
$\alpha=1/\rho_-<1$; if $k\geq\mu$, then $t/s\leq1$.
On $\mu/2\leq k\leq2\mu$ this gives
\[
 \frac{w_{m,k}}T\leq C\mu^{-1/2}
 e^{-c(k-\mu)^2/\mu}.
\]
Its sum outside $|k-\mu|\leq M\sqrt V$ is bounded by a Gaussian
tail tending to zero as $M\to\infty$. The lower indices satisfy
\[
 \sum_{1\leq k<\mu/2}\frac{w_{m,k}}T
 \leq Ce^{-c\mu}\mu^{\alpha-1}
 \sum_{k<\mu/2}k^{1/2-\alpha}
 \leq C\sqrt\mu e^{-c\mu}.
\]
For $2\mu<k\leq\eta_0m$, the deviance is at least $ck$, and
$\sum C\sqrt{k}e^{-ck}/\mu=o(1)$.
For $k=0$, integration over $0\leq s\leq m^{-1/r}$ gives
$J_{m,0}\geq cm^{-1/r}$, hence
$w_{m,0}/T\leq C\mu^{1/r-1}e^{-c\mu}=o(1)$.

For $k>\eta_0m$, fix $s_0>0$ and integrate the compact moment on an
interval of length $c/m$ next to $s_0$. Uniformly for $0\leq k\leq dm$,
\[
 J_{m,k}\geq cm^{-1}s_0^ke^{-m\psi(s_0)}.
\]
Eventually $t<s_0$, and therefore
\[
 \sum_{k>\eta_0m}w_{m,k}(t)
 \leq Cm e^{m\psi(s_0)}
 \frac{(t/s_0)^{\eta_0m}}{1-t/s_0}.
\]
After division by $T$ this tends to zero: $\log(s_0/t)\to\infty$
and its negative exponent dominates $m\psi(s_0)$ and the polynomial
factors. In particular no factor $m$ multiplies an error $e^{-c\mu}$.
This matters when $\mu$ grows slowly.

The central Gaussian Riemann sums \eqref{eq:central} and these tails
give \eqref{eq:fullkernel}.
For $t$ in a fixed compact subinterval of $(0,\infty)$, the relevant
saddles stay in a compact interval, $V\asymp m$, and the same argument
gives central Gaussian sums and tails with $m$ in place of $\mu$.
Indices separated by a fixed amount in $k/m$ from $H(t)$ are handled
without endpoint saddle asymptotics. Integrate their moments on a
width-$1/m$ interval at $s_-=te^{-h}$ or $s_+=te^h$; a fixed small
$h>0$ gives an exponent gain at least $cm$. Their individual evaluation
weights are at most $Cm e^{-cm}$, and there are only $dm+1$ indices.
This proves the fixed-positive-$t$ assertion.
\end{proof}

\section{Proofs on the compact base}
An affine section of $\OO_{\CP^1}(dm)$ is a polynomial of degree at
most $dm$. Indeed it is entire, and holomorphicity of
$z^{-dm}f(z)$ at infinity kills all larger Taylor coefficients.
Apply this observation successively to each variable. The full product
space has basis $z^ke^m$, $0\leq k_j\leq d_jm$.
Membership in $\mm_x^q$ is exactly the condition that all Taylor
coefficients of total degree below $q$ vanish. Its selected basis is thus
\begin{equation}\label{eq:basis}
 z^ke^m,\qquad 0\leq k_j\leq d_jm,\quad \sum_jk_j\geq q.
\end{equation}
Angular integration makes these global sections orthogonal with squared
norm $\prod_jJ_{j,m,k_j}$. Every jet of total degree below $q$ is
globally realized once $q\leq\min_j(d_jm)+1$. Counting these jets gives
\eqref{eq:dimension}; their number is $\binom{q+n-1}{n}=O(q^n)$.
This proves global section production and determines the Hilbert norms
on the full compact manifold.

Let $\mathsf K_j$ be independent integer variables with laws
$\Pr(\mathsf K_j=k)=w_{j,m,k}(t_j)/K_{j,m}(t_j)$.
The exact global kernel identity is
\begin{equation}\label{eq:probability}
 B_m(z)=\left(\prod_jK_{j,m}(t_j)\right)
 \Pr\left\{\sum_j\mathsf K_j\geq q\right\}.
\end{equation}
At the observations of Theorem~\ref{thm:base},
$\mu_j\to\infty$, $\mu_j=o(m)$ and $mG_j(t_j)\sim r_j\alpha_jq$.
Lemma~\ref{lem:moments} gives the product amplitude in
\eqref{eq:transition}. It also gives the Gaussian law of the normalized
sum: on bounded product intervals the central estimates are Riemann
sums for the product Gaussian density; the tails just proved make the
complement arbitrarily small. A half-space boundary has zero Gaussian
measure, so enclosing it in a strip of small Gaussian mass proves
\[
 \frac{\sum_j\mathsf K_j-\sum_j\mu_j}{\sqrt V}
 \Longrightarrow N(0,1).
\]
A linear combination of independent standard Gaussians whose squared
coefficients sum to one is standard Gaussian, by an orthogonal change
of variables in $e^{-|v|^2/2}$. The moving threshold tends to $-\tau$,
and its lattice spacing tends to zero. Equation~\eqref{eq:probability}
now proves the local assertion.

On \eqref{eq:homogeneous}, expansion \eqref{eq:expansion} gives
\[
 \sum_j\mu_j=T_m+\Delta_m,\qquad
 \Delta_m=\sum_j(r_j+1)\varepsilon_j mt_j^{r_j+1}
              +O\left(\sum_jmt_j^{2r_j}\right),\qquad V\sim Rq.
\]
Since $r_j\geq2$, each remainder is little-o of its own positive
correction. Hence
\[
 \frac{\Delta_m}{\sqrt{Rq}}=(1+o(1))\sum_j
 \frac{(r_j+1)\varepsilon_j(\alpha_j/r_j)^{1+1/r_j}}{\sqrt R}
 q^{1/2+1/r_j}m^{-1/r_j}.
\]
This proves the three ranges in Proposition~\ref{prop:shift}.
In the supercritical case tightness of the normalized sum makes the
selected probability tend to one. Also $a(t)\to1$ and
$g_j(t_j)/(r_j^2t_j^{r_j-1})\to1$ in every case.
If the growth ratio fails to tend to zero, pass to a subsequence with
a positive finite or infinite limit; positivity of the displayed constants
then disproves the unshifted law, for example at $\tau=0$.

For the measure assertion, in the affine chart
\[
 \omega=\sum_j\frac{g_j(t_j)}\pi\dA_j,\qquad
 \frac{\omega^n}{n!}=\prod_j\frac{g_j(t_j)}\pi\dA_j,
 \qquad \int_Y\omega^n/n!=D.
\]
The last identity follows from
$\int_0^\infty g_j(t)\,dt=H_j(\infty)-H_j(0)=d_j$.
The full kernel is $\prod_jK_{j,m}$. Lemma~\ref{lem:moments} gives
almost-everywhere convergence of its $m^{-n}$ density relative to
$\nu_Y$ to that of $\omega^n/n!$; the omitted coordinate divisors
have zero smooth-volume measure. Its mass is
$\prod_j(d_jm+1)/m^n\to D$. If $f_m,f$ are these nonnegative
densities, then $\min(f_m,f)\to f$ almost everywhere and is dominated
by $f$. Therefore
\[
 \int|f_m-f|\,\nu_Y
 =\int f_m\nu_Y+\int f\nu_Y-2\int\min(f_m,f)\nu_Y\longrightarrow0.
\]
Finally the full-minus-selected kernel is nonnegative and has integral
$\binom{q+n-1}{n}=o(m^n)$ by orthogonal decomposition. This proves
\eqref{eq:measure}.

\section{Non-descending selectors and exact fibration transfer}
Let $\pi:X\to Y$ be a surjective holomorphic submersion of compact
complex manifolds with connected fibers of dimension $f\geq1$.
Write $N=n+f$ and $E=\pi^*L$, $h_E=\pi^*h$. Fix a Hermitian
form $\omega_X$ on $X$ and a smooth positive volume $\nu_X$ satisfying
\begin{equation}\label{eq:matched}
 \pi_*\nu_X=\nu_Y.
\end{equation}
Choose $p\in\pi^{-1}(x)$ and adapted holomorphic coordinates $(z,w)$
centered at $p$. Let $\chi_p$ equal $\log(|z|^2+|w|^2)$ on a
small coordinate ball, be smooth away from $p$, and be cut off to zero
outside a larger ball. Define
\begin{equation}\label{eq:selector}
 \varphi_m=\frac{q_m+N-1}{m}\chi_p,\qquad
 \II_m=\II(m\varphi_m).
\end{equation}
Here $\II(\Psi)$ uses $|F|^2e^{-\Psi}$ local integrability, with
no extra factor two in the exponent. The selector chooses the section
space; the Hilbert norm below remains smooth.

\begin{theorem}\label{thm:fibration}
For all large $m$, $\varphi_m$ is $\omega_X$-plurisubharmonic and
$\II_m=\mm_p^{q_m}$. Both its pole and its ideal fail to descend.
Nevertheless pullback gives an exact isometric identification
\[
 H^0(X,E^m\otimes\II_m)=\pi^*V_m,
 \qquad \|\pi^*s\|_{h_E^m,\nu_X}=\|s\|_{h^m,\nu_Y}.
\]
Consequently $B_{X,m}=B_m\circ\pi$ and all local laws above transfer
to every upstairs sequence projecting to an admissible base observation.
The section-growth exponent is $n$, rather than $N$, and
\begin{equation}\label{eq:lifted}
 m^{-n}B_{X,m}\nu_X\longrightarrow
 \left(\frac{\omega^n/n!}{\nu_Y}\circ\pi\right)\nu_X
 \quad\text{in total variation}.
\end{equation}
Division by $\dim V_m$ instead gives the probability measure obtained
by dividing the limit by $D$.
\end{theorem}

\begin{proof}
The cutoff logarithm is plurisubharmonic near its sole pole; on the
cutoff annulus its negative Hessian is bounded. Thus
$\ddc\chi_p\geq-C\omega_X$ globally. Since
$(q_m+N-1)/m\to0$, \eqref{eq:selector} is $\omega_X$-psh eventually.

A nonzero holomorphic germ in $N$ variables has a first homogeneous
Taylor term of degree $\ell$. Its squared modulus averaged on a sphere
of radius $R$ is $cR^{2\ell}(1+o(1))$, $c>0$; the initial polynomial
has nonzero squared spherical integral and the higher terms are uniformly
smaller. Its weighted radial integral is finite precisely when
\[
 \int_0^{R_0}R^{2\ell}R^{-2(q+N-1)}R^{2N-1}\,dR
 =\int_0^{R_0}R^{2\ell-2q+1}\,dR<\infty,
\]
or $\ell\geq q$. Smooth positive volume factors do not change this.
Away from $p$ the weight is smooth. Hence $\II_m=\mm_p^q$, a
coherent ideal generated locally by degree-$q$ monomials.
The selector is infinite at $p$ and finite at other points of the same
positive-dimensional fiber. A pulled-back ideal would impose the same
base constraint at each point of that fiber, whereas this ideal is the
unit ideal except at $p$. Neither descends.

Every section of $\pi^*L^m$ descends. Over a trivializing base open
set its coefficient is holomorphic and constant on each connected compact
fiber, by the maximum modulus principle. Local holomorphic sections of
the submersion show that the resulting base coefficient is holomorphic;
fiberwise constancy makes these coefficients glue. In adapted coordinates
the pulled-back coefficient is $f(z)$, independent of $w$, so its order
at $p$ equals its order at $x$. This proves the selected-space equality.
Condition~\eqref{eq:matched} gives the norm equality by integration
along fibers. Pulling back an orthonormal basis proves the kernel identity.

If $b_m,b$ are the base densities in \eqref{eq:measure}, the corresponding
upstairs densities are $b_m\circ\pi,b\circ\pi$. Their $L^1$ difference
is exactly $\int_Y|b_m-b|\nu_Y$ by \eqref{eq:matched}, proving
\eqref{eq:lifted}. Dimension equality and \eqref{eq:dimension} give
the intrinsic exponent and the final probability normalization.
\end{proof}

The matched-volume hypothesis is nonempty for every such submersion.
Start from any smooth positive volume $\mu_X$. Its pushforward is
$b\nu_Y$ with $b$ smooth and positive: in finitely many submersion
boxes, a partition of unity compactly supported in the boxes reduces
fiber integration to integrals of smooth compactly supported densities
in $w$. Differentiation under these integrals proves smoothness, and
nonempty fibers give positivity. Set
\[
 \nu_X=\mu_X/(b\circ\pi).
\]
Then \eqref{eq:matched} holds. In particular $X=Y\times F$ for any
compact connected Hermitian $F$ is a realization; $F$ may be non-K\"ahler.
Product volumes with fiber mass one are matched. Nonproduct matched
volumes include $W(y,u)\nu_Y\nu_F$ with $W>0$ and
$\int_FW(y,u)\nu_F=1$ for all $y$.

For an arbitrary smooth positive volume $\nu'_X$, the exact identity is
\[
 \|\pi^*s\|^2_{\nu'_X}=\int_Y|s|_{h^m}^2\,\pi_*\nu'_X.
\]
Thus the transformed base volume, rather than the original product volume,
governs the kernel. The numerical local constants proved here require
the stated matched volume; the displayed identity alone does not give a
new local theorem for arbitrary transformed volumes.

\section{Scope}
The base theorem treats the specified separable radial compact metrics,
their genuinely nonzero smooth remainders, ordinary total-order point
ideals, and interior energy allocations. It proves global comparison
through exact global monomials and uniform compact moments. Mixed or
angular remainders, other bases and observation faces require further
analysis. Its homogeneous-center range is sharp within this class.

The fibration theorem requires pullback bundles and metrics, compact
connected fibers, and a matched smooth norm. The isolated moving selector
gives an explicit non-descending singularity with exactly computable
selected sections. Arbitrary fiber selectors can impose different base
constraints; singular Hilbert weights and disconnected fibers are outside
the theorem. The lifted limit retains the chosen conditional fiber volume
and is not an ambient-dimensional Monge--Amp\`ere measure. These explicit
sufficient conditions do not close the unrestricted higher-dimensional
programme or classify arbitrary fiber-selected singularities.
\end{document}
